\begin{frame}{Conclusion}
Le livre numérique apparaît plutôt comme une nouvelle offre, qui n'est pas nécessairement à prendre comme une offre concurrente, mais comme une offre supplémentaire.\\ 
La remédiation s'avère être une stratégie fondamentale des nouveaux media numériques.\\ 
Du point de vue des usages (design), elle permet aux utilisateurs d'assurer une continuité formelle et pratique, facilitant la prise en main des nouveaux outils, plateformes, des oeuvres elles-mêmes.\\ 
Du point de vue esthétique, elle ouvre aussi à des hybridations inattendues. Cependant, à l'heure actuelle, les enjeux techniques de la remédiation demeurent dans les mains des grandes firmes technologiques, qui développent autant qu'elles verrouillent les outils de développement, d'écriture et de lecture.
\end{frame}